\cleardoublepage

\section{GCF \& LCM}

Sometimes two numbers meet in the middle. The Greatest Common Factor (GCF) is 
the biggest number they both share as a factor. The Least Common Multiple (LCM) 
is the smallest number they both share as a multiple. These ideas sound opposite, 
but they both help us solve problems with shared values.

Take $12$ and $18$. Their GCF is $6$ since it is the biggest number that divides 
them both. Their LCM is $36$ which is the smallest number they both fit into. We 
find these using prime factorization. Once you learn how, you see how numbers 
connect. And once you see the connection, math becomes clearer.

\subsection{Example: GCF and LCM}

When working with two or more numbers, it's often useful to find numbers they 
have in common. The \textbf{Greatest Common Factor (GCF)} is the largest number 
that divides two numbers exactly, while the \textbf{Least Common Multiple (LCM)} 
is the smallest number that both numbers divide into evenly.

To find the GCF and LCM, you can use prime factorization. For example, let's 
find the GCF and LCM of $12$ and $18$. The prime factorization of $12$ is $2^2 
\times 3$, and for $18$ it's $2 \times 3^2$. The GCF is the product of the 
lowest powers of common primes: $2^1 \times 3^1 = 6$. The LCM is the product 
of the highest powers: $2^2 \times 3^2 = 36$.

Let's apply this to a few more examples. The GCF of $24$ and $36$ is $12$, 
since $12$ is the largest number that divides both. The LCM of $8$ and $12$ 
is $24$, which is the smallest number both can divide into. For $14$ and $28$, 
we see that $14 = 2 \times 7$ and $28 = 2^2 \times 7$, so the GCF is $14$ and 
the LCM is $28$. Mastering GCF and LCM will help you with fractions, ratios, 
and solving many real-world problems.

\newpage
\subsection{Hand-drawing 1}
This is where you will see a hand-drawn answer to the previous question in \textbf{English}.
\begin{figure}[htbp]
  \centering
  \includegraphics[width=\linewidth]{example-image} % TODO: Update caption for actual content
  \caption{A hand-drawn answer in English.}
\end{figure}

\newpage
\subsection{Hand-drawing 2}
This is where you will see a hand-drawn answer to the previous question in \textbf{Korean}.
\begin{figure}[htbp]
  \centering
  \includegraphics[width=\linewidth]{example-image} % TODO: Update caption for actual content
  \caption{A hand-drawn answer in Korean.}
\end{figure}

\cleardoublepage
\subsection{Short Question (English)}
What is the GCF of 8 and 12?

\fullpageimage{images/q1_6_1-eng.jpg}

\newpage
\subsection{Short Question (Korean)}
What is the GCF of 8 and 12?

\fullpageimage{images/q1_6_1-kor.jpg}

\newpage
\subsection{Short Question (Answer)}
What is the GCF of 8 and 12?

\textbf{Answer:} \textbf{4} — the largest number that divides both.

\cleardoublepage
\subsection{Long Question (English)}
Find the LCM of 6 and 10 using prime factorization. Show your steps.

\fullpageimage{images/q1_6_2-eng.jpg}

\newpage
\subsection{Long Question (Korean)}
Find the LCM of 6 and 10 using prime factorization. Show your steps.

\fullpageimage{images/q1_6_2-kor.jpg}

\newpage
\subsection{Long Question (Answer)}
Find the LCM of 6 and 10 using prime factorization. Show your steps.

\textbf{Answer:}\\
- 6 = 2 × 3\\
- 10 = 2 × 5\\
Take the highest powers: 2 × 3 × 5 = \textbf{30}

\cleardoublepage
\subsection*{Challenge Question (English)}
The GCF of two numbers is 7. One number is 35. What are possible values of the other number?

\fullpageimage{images/q1_6_3-eng.jpg}

\newpage
\subsection*{Challenge Question (Korean)}
The GCF of two numbers is 7. One number is 35. What are possible values of the other number?

\fullpageimage{images/q1_6_3-kor.jpg}

\newpage
\subsection*{Challenge Question (Answer)}
The GCF of two numbers is 7. One number is 35. What are possible values of the other number?

\textbf{Answer:}\\
The other number must be divisible by 7 but share no higher factor with 35.\\
\textbf{Examples:} 7, 14, 21, 28, 49, etc.
